\begin{abstract}
Attention weights are often used as a convenient explanation of which words matter most in a Transformer prediction, but prior work has questioned whether attention reliably reflects token importance. We revisit this question in a controlled setting by fine-tuning DistilBERT on the SST-2 sentiment classification task and repeating the experiment across four random seeds. For each validation example, we rank tokens using three signals: last-layer \texttt{[CLS]} attention, gradient $\times$ input, and leave-one-out (LOO) deletion. We compare these rankings with Spearman correlation and evaluate them with token-deletion tests on the full validation set, measuring confidence drops and prediction flips after removing top-ranked tokens. Across seeds, the classifier achieves stable validation performance, with mean accuracy $0.9074 \pm 0.0052$. However, agreement between ranking methods is weak: attention has mean Spearman correlation $0.224$ with gradient $\times$ input and $0.176$ with LOO. Deletion experiments show a clearer hierarchy. Removing the top LOO-ranked token yields the largest mean confidence drop ($0.0383$) and highest flip rate ($0.302$), while attention-based deletion is stronger than random deletion but substantially weaker than LOO. In our setting, attention-based deletion also exceeds gradient $\times$ input deletion, suggesting that raw attention contains some useful signal even though it does not align closely with alternative importance definitions. Overall, our findings suggest that attention is informative but not a reliable standalone proxy for token importance. The project provides a compact and reproducible empirical benchmark for testing interpretability claims in Transformer-based text classification.
\end{abstract}
